\begin{table}[H]
\caption{Newcomer Bucket Robustness ($<$ 1 Week)}
\label{tab:newcomer_robustness}
\centering
\footnotesize
\begin{tabular}{@{}lrrrrr@{}}
\toprule
\textbf{Specification} & \textbf{Arrival HR} & \textbf{95\% CI} & \textbf{Summed HR (upper bnd)} & \textbf{N} & \textbf{Events} \\
\midrule
ModelA\_Newcomer\_Baseline (< 1 Week) & 1.059 & [1.043, 1.075] & 1.107 [1.09, 1.12] & 4,029,008 & 475,633 \\
ModelA\_Newcomer\_Observable (< 1 Week) & 1.102 & [1.089, 1.115] & 1.217 [1.20, 1.23] & 4,029,008 & 475,633 \\
ModelB\_Newcomer\_Baseline (< 1 Week) & 0.987 & [0.975, 0.999] & 0.976 [0.97, 0.99] & 4,029,008 & 475,633 \\
ModelB\_Newcomer\_LengthOnly (< 1 Week) & 0.980 & [0.968, 0.992] & 0.961 [0.95, 0.97] & 4,029,008 & 475,633 \\
ModelB\_Newcomer\_ScoreOnly (< 1 Week) & 0.960 & [0.948, 0.972] & 0.923 [0.91, 0.93] & 4,029,008 & 475,633 \\
ModelB\_Newcomer\_Quality (< 1 Week) & 0.935 & [0.924, 0.947] & 0.877 [0.87, 0.89] & 4,029,008 & 475,633 \\
\bottomrule
\end{tabular}
\vspace{0.35em}
\begin{minipage}{\linewidth}
\footnotesize
\raggedright
Primary HR is the answer-arrival increment ($\beta_4$); the summed DiD ($\exp(\beta_2+\beta_4)$) is shown as an upper bound.\\
Cox-table standard errors are model-based; matched-pair bootstrap uncertainty is reported separately when generated.\\
Events = helping events in the full analysis sample (after time rounding). When interval rows exceed the fit cap, estimation uses a matched-pair subsample, but Events still refer to the full sample.
\end{minipage}
\end{table}